\subsection{(\S13) Difficulties in the recognition of psychological type}

\subsubsection{Joseph Wheelwright}

\begin{enumerate}

\item \ovalbox{\it Requirement of precise use of typology} To use the theory
with precision, one has to be able not only to (1) recognize and accurately
name the main ‘functions’ that a person is using to express his or her
consciousness (thinking, feeling, intuition, and sensation), but also to (2)
figure out which of the two functions that are likely being most often used is
primary and which secondary—and beyond that to (3) make clear the ‘attitude’
with which each function is being deployed. (The choices Jung and Myers give us
here are only two, extraverted and introverted, and it is Myers’s view, and my
own, that if the primary function is extraverted then the auxiliary will be
introverted, and vice versa. This natural alternation of extraversion and
introversion in our functions of consciousness is, by my observation, very
adaptive: it keeps us from becoming too one-sided.)

\item \ovalbox{\it It is a frequent mistake to confuse the first two functions}
This kind of conflation of the two leading functions into one—comprising the
attitude of the dominant and the function of the auxiliary—is a very easy
mistake to make in attempting type diagnosis.

\end{enumerate}

\subsubsection{John Beebe}

\begin{enumerate}

\item \ovalbox{\it One makes their own judgement on who they are}
In MBTI certification training, practitioners are taught that it is unethical
to assert what someone’s type is in contradiction to what the person feels it
is. It is a core feeling value in MBTI facilitation that each person should be
the final judge of who they fundamentally are, including what their type is.
The strength of this tradition in MBTI may owe something to Isabel Briggs
Myers’ own dominant function—introverted feeling.

\end{enumerate}

\subsubsection{Value of the model}

\begin{enumerate}

\item \ovalbox{\it What is ``function''} “The exercise and performance of a
function is something to enjoy, as a pleasant or healthy activity, as the
operation of one’s powers in any sphere of action” (Hillman, 1971/1998, p. 91).

\item \ovalbox{\it Function is used everywhere even between author and reader}
But I am not performing this exercise in a vacuum, just for myself. It is my
way of teaching, of conveying, even of trying to take care of the reader, who I
imagine is reading this chapter in hopes of understanding how to use type in an
analytic or type assessment practice. My auxiliary function, introverted
thinking, is actually trying to take care of you as you read, by getting you to
draw the stick figure of me so that you can visualize both the type theory and
the man who is explaining it to you in the typological terms I have found most
helpful. Whether you feel taken care of, and whether your unclarity about type
is actually lessened by this instruction, is of course a result of how you
receive me, which depends in part on your own typology.

\end{enumerate}

\subsubsection{Parental energy of the auxiliary function}

\begin{enumerate}

\item \ovalbox{\it Auxiliary Fi and Fe} Wordless expression is a common feature
of introverted feeling. It is more interested in solving an underlying feeling
problem than in making communications to address people’s expressions of
discomfort triggered by the problem. Extraverted feeling, by contrast, tends to
be more communicative, even chatty, using a stream of reassuring words to put
people at ease.

Another marker of introverted feeling is that, like any introverted function,
it tends to be more thoroughly original and thus to appear ‘quirkier’ than its
extraverted sibling. $[\ldots]$ Extraverted feeling can be charming, but not
usually, I would assert, through originality and rough-edged sincerity.

Had Jo's feeling been mainly of the extraverted variety, I would assert that he
probably would not have taken the risk of asking the young man the name of his
analyst since this violated a cultural norm of our Institute and could have
elicited an angry refusal. If Jo had been focused on extraverted feeling, I
think he also would have taken more care to keep the young man from drawing the
conclusion that Jo considered some Institute analysts to be unsuitable or
unhelpful—a view that, if it became known, would certainly not have fostered
democratic harmony within the community of analysts.

\end{enumerate}

\subsubsection{Heroic energy of the dominant function}

\begin{enumerate}

\item \ovalbox{\it How dominant function differs from auxiliary function} In
contrast to the auxiliary function, the dominant or ‘superior’ function is less
involved in the care of others and more in the assertion of self. Jung’s
discussion of the archetypal motif of the night-sea journey of the hero in
Symbols of Transformation (1952/1967, ¶¶308–312) provides a beautiful model for
the way a heroic introverted intuition approaches the problem of relating to
the unconscious. In asserting this vision of a consciously irrational ego (i.e.
an intuitive type), he sacrificed his caretaking role with regard to Freudian
psychology, which was the very basis on which Freud had anointed him his ‘crown
prince.’ He was ever after accused by Freud of abandoning the scientific study
of the unconscious, which for Freud could only be accomplished rationally, that
is (to use Jung’s later terminology) through a dialectic of thinking and
feeling.

\item \ovalbox{\it Jung might be Ni-Te} One could say that in response to this
dream Jung’s identity emerged, and that his identity was expressed through a
rather characteristic burst of introverted intuition. In contrast to the type
designation of introverted thinking given Jung by many Jungians, including
occasionally Jung himself, I myself read Jung as an introverted intuitive type,
with extraverted thinking his auxiliary function. What is important here is to
note that, in the way he approaches the dream he tried to share with Freud,
Jung asserts his intuition at the moment he sees it as presenting his ‘own’
standpoint. There is narcissism in this as well as a certain heroic
combativeness. He is emphatically not taking care of Freud, as he does in the
more ‘rational’ writings he was publishing at the time of the dream, where he
uses extraverted thinking to argue for the validity of psychoanalysis. Jung’s
dream, and the way he interprets it, is compensatory to letting himself be used
in this way. The dream fosters the emergence of an extreme self-assertion
carried by the intuitive function, which (at that moment of insight at least)
became the superior function for Jung.

\item \ovalbox{\it How to tell dominant from auxiliary in therapeutic case} With
a patient in the therapeutic situation, we often have to distinguish between
the way the patient asserts self and the way the patient takes care of another.
This is not so hard to do because in the analytic situation the ‘other’ will
usually be the analyst. There is something heroic about self-assertion (it
should be noted that many patients have a considerable difficulty asserting
themselves with anything like the definiteness Jung describes in taking
ownership of his dream), and there is something parental about taking care of
the other. The analyst may want to note the ways the patient is parental in the
transference, and not just the ways the patient is infantile. (Developing
Winnicott’s (1987) notion of the analyst’s way of ‘holding’ the patient
throughout an analysis, contemporary relational psychoanalysts have implied
that the patient holds the analyst during treatment quite as much the analyst
holds the patient [Samuels, 2008, personal communication] and of course there
are wide variations both in patients’ capacities to do this and in the ways
they do this, into which both the strength and the typology of the patient’s
auxiliary function figures.) This distinction helps us to differentiate the
patient’s (heroic) dominant function from his or her (parental, caretaking)
auxiliary function, and that can be an enormous help in establishing a reliable
type diagnosis.

\end{enumerate}

\subsubsection{Recognizing the eight function-attitudes}

\begin{enumerate}

\item \ovalbox{\it We don't yet have a systematic way to recognize them} One
has, in other words, to be able to recognize, and tell the difference, between
introverted thinking, introverted feeling, introverted sensation, introverted
intuition, and extraverted thinking, extraverted feeling, extraverted
sensation, and extraverted intuition. Learning to do this requires conscious
practice. It is not unlike the way one learns how to read music. Unfortunately,
we don’t have a mnemonic song like “Do-Re-Mi,” which Mary Martin, and later
Julie Andrews, sang in The Sound of Music, to learn type recognition the way we
learn to recognize the basic tones of the Western musical scale.

\end{enumerate}

\subsubsection{Beyond dominant and auxiliary}

\begin{enumerate}

\item \ovalbox{\it Beyond dominant and auxiliar is harder}
The functions are not as easily recognized in therapy. A real-life person,
unlike a stereotyped character identified with a single function, has access to
all eight functions of consciousness, even if some are in shadow, and will
deploy one or another depending on the context and the type of consciousness
called for by that context. Also, the patient in analysis is often in the grip
of complexes, which notoriously produce what Jung, quoting Janet, called an
abaissement du niveau mentale, a reduction of the mental level, such that the
energy that normally attaches itself to the superior and auxiliary function,
allowing them to surface, is absent. When these functions are not active, the
tertiary and inferior functions emerge. ‘Tertiary’ and ‘inferior’ are terms
that imply that there is a gradient of differentiation in the four types of
consciousness that normally describe someone’s ‘ego,’ at least as that
sometimes inexactly defined term is understood by analytical psychology. (I
prefer to think of them as less conscious parts of a ‘little-s’ self.) As the
least differentiated of the functions consciously available to the patient, the
tertiary and inferior functions tend to be less adapted to reality and more
influenced by unconscious complexes, which are in fact usually dominating the
psyche when the third and fourth functions emerge in recognizable form. Their
presentation is thus often floridly neurotic, easily characterized as
obsessive, cyclothymic, hysterical, or paranoid, creating an obvious link to
psychopathology. When it is easy to diagnose neurotic traits or character
pathology in an analytic patient, it is a tip-off that one is looking not at
the patient’s natural (superior and auxiliary function) typology but at “a
falsification of the original personality” (Jung 1950/1959, ¶214). Naturally, a
person can also falsify his or her original personality in a more adaptive way
by conforming to the expectations of a family, school, occupation, or culture.

\item \ovalbox{\it Precaution in making type diagnosis} We should be cautious,
therefore, in making any type diagnosis. It is best not to try to type someone
who has not yet made a connection to the self that would be natural to him or
her, because all you may be doing is noting the “negative personality” (Jung,
1950/1959, ¶214) that has swallowed up the patient’s true self (Beebe, 1988).
Sometimes, however, knowing that the inferior and tertiary functions reflect
what someone ‘is least good at’ can be a clue to the actual type. The person
who constantly obsesses about small feeling matters, finding other people’s
feelings an endless burden, may be not an extraverted feeling type, for whom
other people’s feelings naturally matter and are thus relatively easy to deal
with, but someone with inferior extraverted feeling, that is, an introverted
thinking type, who is in constant danger of ignoring the feelings of others.
Marie-Louise von Franz (1971/1998) has written the definitive text on the
inferior function, and in many ways her monograph is also the best book on type
for clinicians because it portrays the way many kinds of patients present
themselves in the office when in the grip of the inferior function. It should
be required reading in all Jungian training programs.

\end{enumerate}

\subsubsection{Connection to the Self}

\begin{enumerate}

\item \ovalbox{\it Do not rush} To discover a patient’s typology, it is better
to wait until the patient shows an original gift for accurately construing or
managing some aspect of what comes up in therapy, rather than attempt to ‘type’
the person when he or she is manifesting a collective persona that could belong
to anybody in the patient’s situation, or when the patient is so evidently
suffering from psychopathology that a syndrome has all but replaced the person.

\item \ovalbox{\it The right moment} The typology of the true self (defined as
the personal, ‘little-s’ self in touch with the transpersonal ‘big-S’ Self
(Gordon, 1985; Beebe, 1988) ) is rarely so stereotyped; rather it opens up the
use of the most differentiated parts of the personality in an individual way
that is a revelation and a pleasure to experience. It’s when the patient is
exhibiting his or her strengths as an authentic person that we can begin to
appreciate the skill with which feeling, thinking, sensation and intuition are
being used. At such moments, we can also see what effective extraversion and
introversion are like when they are used as conscious attitudes.

\end{enumerate}

\subsubsection{Extraversion and introversion}

\begin{enumerate}

\item \ovalbox{\it How to tell extraversion} When the patient is using a
well-differentiated extraverted function, the function will seek to merge with
some aspect of the analyst in a way that the analyst does not find particularly
uncomfortable. When extraverted feeling is differentiated, the analyst feels
appreciated and respected, and there is a sense of one’s good will being seen
and met. When extraverted thinking is highly differentiated, the analyst will
find that it is safe to let the patient set the agenda, like a general
directing the campaign of the therapy. When the patient’s extraverted sensation
is well developed, the analyst has the experience of a ready participation in
what is happening in the moment and an accompanying impatience with
abstractions, as if what is already there is sufficient without much
interpretation. Extraverted intuition can feel intrusive, but it is also
entertaining and astonishing in the way it can pick up on fresh possibilities
for developing the objectives of the therapy in the world.

\item \ovalbox{\it How to tell introversion (more difficult)} Introversion,
when used consciously, is not as easy to discriminate, and indeed the functions
of introverted feeling, introverted sensation, and introverted intuition are
easily confused with each other. Introverted thinking can usually be
distinguished by the fact that it tends not to know when to stop, and needs to
define everything freshly, to the point that it becomes exhausting and hard to
follow. It, like the other introverted functions, seeks to match its experience
of an object with an a priori, archetypal understanding of that category of
object already present in the unconscious. The introverted move away from the
outer object is therefore the first step in a process that takes the
introverted function’s libido deep into the unconscious of the introverted
subject, to see if the object really matches up. (That it often doesn’t live up
to the archetype helps to explain the frequent disappointment that introverted
functions register in analysis, a disappointment that must not be confused with
a condition to be treated, even if it is dysphoric to the subject. People with
superior introverted functions must register this disappointment when the
object simply doesn’t match up. It’s their normal way of reacting.) Introverted
intuition seeks to match up the experience with an image of an archetype,
something like a visual metaphor. Introverted sensation likes to establish
whether the experience of the object checks out with an inner sense of what has
already been established through long human experience as ‘real.’ And
introverted feeling wants to know if the object as experienced is conducting
itself in accord with what is fitting for such an object, that is, if a bride
is acting like a bride, if a home feels like a home, if the boss is behaving as
she should in her role.

\item \ovalbox{\it Suggestion for recognizing introversion} The practitioner
should grow accustomed to the way introverted functions are constantly sizing
up what happens in a therapeutic or in a type assessment encounter, to see if
it checks out with the rich inner world of already-known, archetypal
experience, against which an introverted function measures everything.
Recognizing the introverted types in their normal functioning is one way to
realize Jung’s enormous contribution to opening up the introverted world as a
part of healthy functioning. Simply not pathologizing introversion is perhaps
the most healing thing a therapist or type counselor can do in our
pathologically extraverted, world-despoiling times. And it is a sign that the
practitioner is well along in the knack of type recognition.

\end{enumerate}
